
The Free Church Is Not a "Church Body"

[Editorial article in "Folkebladet," April 5, 1899.]

The Committee of the United Church for the Revision of the Clergy List makes the assertion that it must strike some ministers from the clergy list of the United Church because they have taken part in the meetings of the Free Church. For, says the Committee: "It must be evident that you cannot stand as a member of two different church bodies at the same time."

This is not "evident." Even if there really were question here of two church bodies, that is surely no sufficient reason why one could not be a member of both, provided that they have the same confession and work for what is good.

So many "church bodies" or parties are, after all, organized here. Suppose now that one wished to work for the salvation of the Jews, another for the heathen, one for mission in China, another for mission in Africa; one for a seminary in Chicago, another for a seminary in St. Paul; one for preaching in English, another for preaching in Norwegian, and so forth—would it then without more ado be "evident" that a man or a congregation cannot stand as a member of two or more of these?

Behind this "evident" thing there must presumably lie the thought that each "church body" is to and must think of itself as being the whole Thing, as we say. For if the meaning is that each "church body" must think of itself as being the whole Church, or the whole Kingdom of God, or at least the Lutheran Church, or at the very least the whole true Church, then of course it is so that it becomes meaningless to be a member of more than the one true body. But we cannot now believe that the United Church has fallen into such a delusion about itself.

If that does not lie behind it, then it must be the meaning one wishes to assert, that a church body is a party which must necessarily combat all other similar parties; for then too it is true that it does not really work to be on both sides in a war of extermination. But we can hardly imagine that either, that the United Church has such an opinion of itself.

But now this must also be added: the Free Church has never been any "church body" in the outdated sense of the word. If it were, then it would presumably have to have a constitution that gathered and delimited it, and it would have to have an authority over the congregations, such as all so-called "church bodies" more or less have.

This is not found in the Free Church. It is a free association of Lutheran congregations without any governing authority over them. The congregations are as independent of one another as England and the United States. They may cooperate when and where they will; but the one cannot compel the other to join where it does not itself see a call to do so.

The meetings of the Free Church are therefore free open meetings for Lutheran Christians. Their resolutions are binding upon the committees which the meetings themselves appoint; for all others they are only to be regarded as encouragements, urgings, proposals, and counsel. Their influence therefore rests exclusively upon the spirit and wisdom that may be found in them.

This can only by great distortion and misuse of words be called a "church body." And it is by no means "evident" why a man could not take part in such a free meeting and nevertheless be fit to stand as a minister in the United Church.

[Georg Sverdrup.]
